% ---
% id: meat-core/frank-wolfe/cor:certified-keylength-reference
% kind: cor
% title: "Key length on a reference companion from the certified bracket"
% summary: >-
%   On a reference companion, the best certified bracket, formed over its
%   enclosed measured family as [[lem:certified-bracket-enclosed]] allows,
%   gives the key length of the protocol whose declaration is the companion,
%   as the first-order bound on the companion does.
% prominence: standard
% status: accepted
% uses: [cor:certified-keylength, thm:firstorder, cor:fo-reference, cor:fo-keylength, lem:certified-bracket, lem:certified-bracket-enclosed, cor:fo-dropped, meat-core/prop:reference-companion, meat-core/cor:reference-companion-factors, meat-core/def:admissible-declaration, meat-core/cor:keylength-bound]
% mentions: []
% assumes: []
% refs: []
% provenance: newly-derived
% lean:
% notes: "Read for the companion of meat-core/prop:reference-companion, which may replace several measured factors, so that the hypotheses of cor:reference-companion-factors are named per replaced factor. Kept apart from cor:certified-keylength so that cor:fo-reference stays out of its closure on routes with no companion."
% ---

Corollary~\ref{cor:certified-keylength} applies to a reference companion as
to any declaration; this corollary carries its conclusion over to the
protocol whose declaration the companion is, as
Corollary~\ref{cor:fo-reference} does for the first-order bound. Where the
companion's measured family is known only through enclosures, the two lower
bounds at a point are formed from enclosures that hold over every family the
enclosures leave possible, and part~(c) of
Lemma~\ref{lem:certified-bracket-enclosed}, or part~(d) under a dropped set,
is what makes them bounds at the companion's own gradient.

\begin{corollary}\label{cor:certified-keylength-reference}
Let $\Decl^{\mathrm{ref}}$ be a compiled reference companion of
Proposition~\ref{prop:reference-companion}, and let $\alpha$, $\acSet$ and
$\Itset$, with its points $x_t$, be as in
Corollary~\ref{cor:certified-keylength} with
$\Decl\coloneqq\Decl^{\mathrm{ref}}$. At each $t\in\Itset$ let
$b^{\mathrm A}_t$ and $b^{\mathrm F}_t$ be real numbers formed in one of two
ways: as in that corollary, with \eqref{eq:lmo-bounds} at the gradient of
$\Decl^{\mathrm{ref}}$ at $x_t$; or, where blocks of the measured family of
$\Decl^{\mathrm{ref}}$ are known through enclosures, entrywise or by factor
enclosures, as parts (ii) and (iii) of Lemma~\ref{lem:certified-bracket} form
them from enclosures $(\hat g,\delta)$ and $(M_p,\rho_p)$ that hold at every
pair of a set $\mathcal P$ as in
Lemma~\ref{lem:certified-bracket-enclosed}\,(c), at a point $x_t$ at which
the hypotheses of that lemma hold and (H-POS) holds for
$\Decl^{\mathrm{ref}}$, the blocks it reads with factor enclosures being read
with those of Corollary~\ref{cor:reference-companion-factors}, under that
corollary's hypotheses (for each replaced factor whose blocks are so read,
the rule of (I1) of the reduction instance of its last link a non-negative
combination, and the blocks of its handed reference carrying those factor
enclosures), so that its conclusion is (H-FAC). Put $\underline{h}^{\mathrm{cert}}_{\alpha}\coloneqq
\max_{t\in\Itset}(b^{\mathrm A}_t+b^{\mathrm F}_t)$, as in \eqref{eq:hcert}.
\begin{enumerate}[label=(\roman*),leftmargin=2.4em]
\item Where $\Decl^{\mathrm{ref}}$ satisfies
  Definition~\ref{def:admissible-declaration} and the hypotheses of
  Corollary~\ref{cor:fo-keylength} hold, the conclusions of
  Corollary~\ref{cor:fo-reference}\,(iii) hold with
  $\underline{h}^{\mathrm{cert}}_{\alpha}$ in place of
  $\hlb[\Decl^{\mathrm{ref}}]$, for the protocol whose declaration is
  $\Decl^{\mathrm{ref}}$.
\item In the setting of Corollary~\ref{cor:fo-dropped}, with a dropped set
  $\mathcal{D}$ and with the numbers formed at the gradient of the dropped
  objective and, in the second way, with part (d) of
  Lemma~\ref{lem:certified-bracket-enclosed} in place of its part (c):
  $\max\{0,\underline{h}^{\mathrm{cert}}_{\alpha}\}\le
  \underline{h}^{\mathcal{D}}_{\alpha}[\Decl^{\mathrm{ref}}]\le
  \hlow[\Decl^{\mathrm{ref}}]$, and where $\Decl^{\mathrm{ref}}$ satisfies
  Definition~\ref{def:admissible-declaration} and the hypotheses of
  Corollary~\ref{cor:keylength-bound} other than \eqref{eq:lower-bound-h}
  hold, the conclusions of
  Corollary~\ref{cor:keylength-bound} hold at
  $h=\underline{h}^{\mathrm{cert}}_{\alpha}$, for the protocol whose
  declaration is $\Decl^{\mathrm{ref}}$.
\end{enumerate}
\end{corollary}

\begin{proof}
\emph{The numbers are bounds at the companion's gradient.} At a point where
$b^{\mathrm A}_t$ and $b^{\mathrm F}_t$ are formed in the second way,
hypothesis (H-FAC) of Lemma~\ref{lem:certified-bracket-enclosed} holds by
Corollary~\ref{cor:reference-companion-factors}, under the hypotheses on
which the reduction and the compiler carry the factor enclosures, so part
(c) of that lemma shows that the enclosures
hold at the gradient of $\Decl^{\mathrm{ref}}$ at $x_t$, (H-POS) holding for
$\Decl^{\mathrm{ref}}$ there, and parts (ii) and (iii) of
Lemma~\ref{lem:certified-bracket}, which $\Decl^{\mathrm{ref}}$ is admitted
to by Proposition~\ref{prop:reference-companion}\,(ii), turn them into
$b^{\mathrm A}_t\le\min_{\mathbf q'\in\acSet}\langle\lmoq,\mathbf q'\rangle$
and $b^{\mathrm F}_t\le\min_{J'\in\Feas^{J}}\langle\lmoJ,J'\rangle$, the
gradient and the exact $\Feas^{J}$ being those of $\Decl^{\mathrm{ref}}$:
that is \eqref{eq:lmo-bounds} at its gradient, as in the first way.

\emph{(i).} Corollary~\ref{cor:fo-reference}\,(iii) holds for any bound
formed as its statement admits, lower bounds on the minima over the
companion's exact sets standing in for the pairings at points of its relaxed
set at which (H-POS), (H-SUPP) and (H-NOPERT) hold for its objects. Each
$x_t$ satisfies the hypotheses of Lemma~\ref{lem:certified-bracket} for
$\Decl^{\mathrm{ref}}$, which gives that by Step~1 of its proof, and the
numbers satisfy \eqref{eq:lmo-bounds} at the gradient of
$\Decl^{\mathrm{ref}}$, so Corollary~\ref{cor:certified-keylength}\,(i)
applied to $\Decl^{\mathrm{ref}}$ shows that
$\underline{h}^{\mathrm{cert}}_{\alpha}$ is such a bound: the value of
\eqref{eq:hlb} at the points of $\Itset$ with both pairings replaced.

\emph{(ii).} The same argument with part (d) of
Lemma~\ref{lem:certified-bracket-enclosed}, which runs part (c) for the
dropped objective, and part (iv) of Lemma~\ref{lem:certified-bracket} gives
the numbers as lower bounds on the two minima at the gradient of the dropped
objective of $\Decl^{\mathrm{ref}}$, and
Corollary~\ref{cor:certified-keylength}\,(iii) applied to
$\Decl^{\mathrm{ref}}$ gives the two inequalities and, under the hypotheses
that corollary's second clause and (ii) here assume, the conclusions of
Corollary~\ref{cor:keylength-bound}.
\end{proof}
